\chapter{Introduction}

\section{Software Engineering}

Software Engineering is a discipline focused on the correct development,
implementation, and maintenance of software systems.

The major pillars of Software Engineering include:

\begin{itemize}
    \item \textbf{Software Life Cycle:} Requirements Engineering,
    Architectural Design, Implementation, Testing/Verification,
    Deployment, and Long-term Maintenance.

    \item \textbf{Quality Assurance:} Formal methods, static analysis,
    and unit/integration testing.

    \item \textbf{Abstraction and Modular Design:} Dividing a system
    into decoupled components with well-defined interfaces and
    abstract data types.
\end{itemize}

\begin{ppnote}
Software Engineering covers the complete development and maintenance
of software systems, with correctness and quality being important
considerations throughout the software life cycle.
\end{ppnote}


\section{Program Analysis}

Given a program $P$, program analysis attempts to determine whether
some property $\phi$ holds for all behaviours of $P$.

This can be written as:

\[
P \models \phi
\]

For example, we may want to verify the property

\[
P \models \text{``no divide by zero''}.
\]

Program analysis can be broadly divided into two modes:

\begin{itemize}
    \item Static Analysis
    \item Dynamic Analysis
\end{itemize}


\subsection{Static Analysis}

Static analysis analyses a program without executing it.

Examples of static analysis techniques include:

\begin{itemize}
    \item Compiler data-flow analysis
    \item Abstract interpretation
    \item Symbolic execution
    \item Model checking
\end{itemize}


\subsection{Dynamic Analysis}

Dynamic analysis obtains information by observing the program during
execution.

Examples include:

\begin{itemize}
    \item Testing
    \item Profiling
    \item Runtime monitoring
    \item Dynamic taint analysis
\end{itemize}

\begin{ppnote}
Static analysis reasons about program behaviour without executing the
program, whereas dynamic analysis uses information obtained from
program executions.
\end{ppnote}


\section{Formal Verification}

Formal Verification can be viewed as:

\[
\text{Formal Verification}
=
\text{Model} + \text{Specification} + \text{Proof}.
\]

\subsection{Model}

The \textbf{model} represents the semantics of the program.

\subsection{Specification}

The \textbf{specification} describes the correctness properties that
the program is expected to satisfy.

\subsection{Proof}

The \textbf{proof} consists of techniques used to establish that the
program, under the given semantics, satisfies the specification.

Examples of formal verification techniques include:

\begin{itemize}
    \item Model checking
    \item Theorem proving
    \item Symbolic verification
\end{itemize}


\section{Program Analysis vs.\ Formal Verification}

Program Analysis may infer different facts about a program.

For example, program analysis may determine that:

\begin{itemize}
    \item An amount influences the balance.
    \item There are two paths.
    \item There is no pointer dereference.
\end{itemize}

Formal Verification instead involves specifying a property and proving
that the property holds.

For example, given

\[
\text{balance} \geq 0
\quad\land\quad
\text{amount} \geq 0,
\]

we may want to prove

\[
\text{balance}' \geq 0.
\]

Thus, the verification goal can be expressed as:

\[
\text{balance} \geq 0
\;\land\;
\text{amount} \geq 0
\quad\Longrightarrow\quad
\text{balance}' \geq 0.
\]

\begin{ppnote}
Program Analysis can infer information about program behaviour,
whereas Formal Verification uses a specification and proof techniques
to establish that a required property holds.
\end{ppnote}


\section{Why are Program Analysis and Formal Verification Hard?}

Program Analysis and Formal Verification have a fundamental tension
between \emph{precision} and \emph{scalability}.

\begin{ppnote}
The more precise we want an analysis to be, the more expensive the
analysis can become.
\end{ppnote}

There are several reasons why Program Analysis and Formal Verification
are difficult:

\subsection{Infinitely Many Inputs}

There may be infinitely many possible inputs to a program. Therefore,
checking every possible input individually is generally not feasible.

\subsection{Large Number of Paths}

The number of possible execution paths can grow exponentially.

\subsection{Loops}

Loops can create infinitely many possible executions because a loop may
execute an unbounded number of times.

\subsection{Complex Program State}

Real programs manipulate complex state involving:

\begin{itemize}
    \item Pointers
    \item Heap objects
    \item Complex data structures
    \item Recursion
    \item Callbacks
    \item Exceptions
    \item Threads
    \item Files
    \item Networks
\end{itemize}

These features make precise analysis and verification challenging.


\section{How Can LLMs Assist?}

Large Language Models (LLMs) can assist Program Analysis and Formal
Verification in several ways.

\subsection{Specification Synthesis}

LLMs can help with specification synthesis, particularly by converting
natural-language requirements into logical specifications.

For example, the requirement

\begin{quote}
Every accepted request must eventually receive a response.
\end{quote}

can be represented as:

\[
G(\text{req}) \rightarrow F(\text{resp}).
\]

Here, $G$ represents that the property should hold globally, while
$F$ represents that the response should eventually occur.

\subsection{Semantic Inference}

LLMs can help with semantic inference in the form of:

\begin{itemize}
    \item Invariants
    \item Preconditions
    \item Postconditions
\end{itemize}

\subsection{Generating Auxiliary Lemmas}

LLMs can also generate auxiliary lemmas that may help verification.

Sometimes an invariant is already known to the underlying decision
procedure, but an SMT solver may still be unable to complete the proof.
In such situations, an LLM can help generate candidate invariants or
auxiliary lemmas.

\begin{ppexample}
Consider a program that computes the sum of an array. An LLM can
generate a candidate invariant that may help the underlying verification
procedure complete the proof.
\end{ppexample}

\begin{ppnote}
LLMs are excellent at generating candidate specifications, invariants,
and lemmas. However, the generated candidates eventually need to be
validated using reliable verification techniques.
\end{ppnote}


\section{What LLMs Are Not Good At?}

LLMs have important limitations when used for Program Analysis and
Formal Verification.

\subsection{LLMs Cannot Verify Their Own Output}

LLMs cannot reliably verify their own output. Therefore, they cannot
be relied upon as proof checkers.

\subsection{Soundness and Completeness}

LLMs are neither sound nor complete.

\begin{itemize}
    \item \textbf{Soundness:} No false positives.
    \item \textbf{Completeness:} No false negatives.
\end{itemize}

Thus, an LLM-generated result cannot automatically be treated as a
correct proof.

\subsection{Precise Reasoning}

LLMs can be weaker at precise reasoning.

For example, an LLM may correctly reason about the following expression
in isolation:

\[
\forall x \exists y.\; R(x,y)
\]

but may accidentally interchange or modify the quantifiers during
subsequent reasoning steps.

\begin{pppitfall}
An LLM-generated proof or specification should not be accepted merely
because it looks plausible. Generated candidates must be checked using
appropriate analysis or verification tools.
\end{pppitfall}


\section{Summary}

LLMs can assist Program Analysis and Formal Verification at four key
layers.

\subsection{Property Discovery}

LLMs can help answer:

\begin{quote}
What should we check?
\end{quote}

This involves discovering properties that should be verified.

\subsection{Abstraction Discovery}

LLMs can help determine:

\begin{quote}
What information should the analyzer maintain?
\end{quote}

\subsection{Proof Search}

LLMs can help identify:

\begin{quote}
What invariants or lemmas make the proof work?
\end{quote}

\subsection{Tool Orchestration}

LLMs can help determine:

\begin{quote}
Which analyzer, verifier, or solver should be used?
\end{quote}

\begin{ppnote}
The overall principle is: \emph{Trust but verify!}

LLMs can be useful for generating candidate properties, abstractions,
invariants, lemmas, and tool choices, but their outputs need to be
validated.
\end{ppnote}


\section{Demo: Doubling Sum in Dafny}

The lecture concludes with a demonstration of proving a
\emph{Doubling Sum} program in Dafny.
